\chapter{Entrevistas semiestructuradas (transcripciones completas)}
\label{apx:entrevistas}

\section*{Nota metodológica}
Para el componente cualitativo del \textit{User Research} se realizaron tres entrevistas virtuales semiestructuradas. Los participantes fueron seleccionados intencionalmente por sus perfiles vinculados con la docencia, la investigación y la gestión académica. Se trabajó con una guía orientativa de preguntas de carácter flexible; durante las conversaciones se modificó el orden de los temas cuando resultó necesario y se realizaron repreguntas para profundizar aspectos relevantes. No se conserva una versión documental verificable que pueda presentarse como la guía definitiva exacta aplicada, por lo que dicha guía no se reconstruye ni se incorpora retrospectivamente en este anexo.

Las condiciones de consentimiento informado se detallan en la sección de Metodología del Capítulo~\ref{ch:user_research}. Las transcripciones completas se reproducen a continuación; para conservar la trazabilidad metodológica al leerlas, ENT-01 y ENT-02 poseen una etapa inicial de exploración previa a la presentación de ImprontIA, mientras que ENT-03 no cuenta con un tramo de \textit{discovery} comparable, dado que la conversación comenzó con una descripción general del proyecto.

\section*{ENT-01 --- Lic. Ariel Vázquez}
\addcontentsline{toc}{section}{ENT-01 --- Lic. Ariel Vázquez}
\begin{ficha}{Ficha de entrevista}
\small
\textbf{Código:} ENT-01\\[2pt]
\textbf{Participante:} Lic. Ariel Vázquez\\[2pt]
\textbf{Rol / perfil:} responsable de UADE Online\\[2pt]
\textbf{Fecha:} 9 de junio de 2026\\[2pt]
\textbf{Modalidad:} virtual\\[2pt]
\end{ficha}
\subsection*{Transcripción completa}
\begin{description}[leftmargin=0cm, labelsep=0.5cm, style=nextline]
\item[\textbf{Entrevistador:}] Sí, nos viene súper bien este tiempito. Vamos a hacer una mini introducción de nuestro trabajo de tesis. La verdad va a ser muy mini en principio, para no sesgar demasiado con la solución, porque queremos abordar la problemática. ¿Te corto un segundo? Nos presentamos porque creo que vos sabés por el mail, pero bueno, me imagino que recibís miles de mails por día.
\item[\textbf{Entrevistado:}] Dale, sí, sí. No, de nada, un placer colaborar.
\item[\textbf{Entrevistador:}] Somos Martina y Juan Ignacio, estudiamos Informática, estamos en el último año y nos quedan pocas materias. Nos gustaría escuchar un poco sobre vos, tu rol en UADE, hace cuánto estás, etcétera.
\item[\textbf{Entrevistado:}] Bien. Yo estoy en UADE desde el año 2012 como profesor y en la gestión desde el año 2018. Siempre trabajé en la Facultad de Ciencias Económicas como director de carrera o de departamento. Desde septiembre del año pasado, hace casi un año, estoy como director de UADE Online. Digamos que soy el que maneja toda el área de educación a distancia de acá de UADE.
\item[\textbf{Entrevistador:}] Perfecto, súper interesante. O sea, de todas las carreras, controlas todo lo que es la plataforma UADE Online y las materias que se dictan.
\item[\textbf{Entrevistado:}] Claro. UADE Online acá es procesada como una sede (como Pinamar, ponele), nosotros lo trabajamos como sede. No es una facultad, es una sede, pero en ella se dictan materias de todo tipo y color, y tenemos materias para todas las facultades. También tenemos oferta para los chicos de la modalidad presencial y, por otro lado, carreras que son 100\% a distancia, donde los alumnos cursan la carrera íntegramente de forma remota. Bueno, en informática está la Licenciatura en Gestión de Tecnología de la Información.
\item[\textbf{Entrevistador:}] Sí, la Licenciatura en Sistemas, habíamos escuchado. Y te hago una pregunta sobre las evaluaciones: yo, por ejemplo, he hecho materias \textit{online} y las evaluaciones fueron presenciales, creo que en todos los casos.
\item[\textbf{Entrevistado:}] Claro, está bien. Para los chicos de la modalidad presencial que cursan materias \textit{online}, el examen es presencial. Solamente evaluamos a distancia a los chicos que eligen la carrera 100\% a distancia. ¿Me explico? Entonces, en esas carreras que son 100\% a distancia, lo que te ofrece UADE es no tener que venir en ningún momento a la universidad.
\item[\textbf{Entrevistador:}] Sí, entiendo que lo elige mucha gente del exterior también.
\item[\textbf{Entrevistado:}] Mira, en realidad no tanto, porque no lo difundimos mucho en el exterior. La realidad es que nuestros alumnos son mayormente de CABA y de la provincia de Buenos Aires, con alguna presencia en el interior, como gente de Bahía Blanca. No tenemos una gran concentración en el interior y menos en el exterior; tenemos casos, sí, claro, pero serán el 4\% o 5\%. No es tanto porque tampoco nos hemos enfocado en comunicar la propuesta en esos mercados. Pero, en general, la propuesta de carreras 100\% a distancia sí suele tener mucha convocatoria del interior y del exterior.
\item[\textbf{Entrevistador:}] Súper interesante, excelente. Bueno, un poco sobre nuestro proyecto, está ligado con evaluar y analizar cómo es la problemática hoy en día sobre la integridad académica en evaluaciones digitales y remotas. Creció mucho el volumen de las evaluaciones de forma remota postpandemia y va creciendo cada día más a medida que se habilitan más instituciones y universidades. Nosotros estamos evaluando qué ventajas tiene, cuáles son sus fuertes, qué problemáticas se presentan y qué les preocupa tanto a la división docente como a los directores. Así queda un poco encaminado por ahí. Más hacia el final de la entrevista te vamos a dar más detalles, pero por ahora es eso. La primera pregunta sería cómo son las evaluaciones virtuales en UADE Online y qué te parecen los desafíos y los aspectos positivos.
\item[\textbf{Entrevistado:}] Bueno, mira, cuando yo tomé el área, las evaluaciones no tenían ningún sistema de control automatizado, sino que se evaluaba al alumno y después se validaba mediante un examen oral; no había ninguna herramienta tecnológica. Justamente lo que están investigando ustedes me gustó porque cuando vine acá, una de las primeras medidas que tomé fue la implementación de un sistema tecnológico para el control de exámenes. Desde marzo de este año lo implementamos. El año pasado me basé en el \textit{benchmark} de todos los sistemas existentes habidos y por haber; en diciembre elegimos la herramienta y en febrero empezó la implementación. Yo considero que, sí o sí, la evaluación remota tiene que contar con una herramienta tecnológica que valide y controle el examen. No hay una manera humana de controlarlo fehacientemente a distancia. Si bien podés tomar un oral de validación, nunca es lo mismo que hacer el 100\% del examen solo. La complejidad de evaluar a distancia sin una herramienta es enorme y, obviamente, atenta contra la integridad y la calidad académica. Es un tema que preocupa mucho y que en general está muy mal manejado. Además de hacer \textit{benchmark} de herramientas, he investigado a otras universidades: muchas no tienen nada implantado y las que tienen usan herramientas que son bastante flojitas.
\item[\textbf{Entrevistador:}] ¿Y esta herramienta que adoptaron qué tipo de informes les da? Es decir, ¿qué tipo de métricas reporta, qué cosas les sirven y qué cosas no tanto?
\item[\textbf{Entrevistado:}] Te cuento de la herramienta, te lo digo también porque no es confidencial: se llama Smowltech (Smowl). Es una herramienta española bastante interesante porque tiene dos modalidades para tomar el examen. Una consiste en el bloqueo de pantalla; te da 5 oportunidades para escaparte de la pantalla y a la sexta te saca del examen. Es una opción un poco violenta. A mí no me gusta tanto y no la implementamos acá porque quedás muy supeditado a los fallos de la tecnología. Por ejemplo, si el alumno dejó las notificaciones abiertas o si no tiene mala intención pero se le suspendió la máquina. La condición de expulsión son cinco salidas o estar 1 minuto afuera de la plataforma (por ejemplo, si hacés Alt + Tab). Una vez que te saca, ya no podés volver a entrar. Para mí, esa modalidad tiene dos desventajas. Por un lado, una desventaja tecnológica en cuanto a que ---ustedes estudian informática y lo saben--- a veces uno no puede prever ciertas contingencias. Te colgaste, dejaste el Teams abierto y cualquier elemento que se ponga por detrás de la pantalla de bloqueo el sistema lo toma como que te quisiste ir, cuando en realidad no estás haciendo nada malo. Por otro lado, hay un factor psicológico: el alumno rinde muy tenso por el miedo a que el sistema lo expulse. Esto es fatal. De hecho, hemos hecho algunos experimentos con esa modalidad y la verdad es que el alumno se pone muy nervioso. Si al condimento propio del examen le agregás la preocupación de que en cualquier momento te puede echar, se vuelve bastante caótico. Uno lo que tiene que apelar es a controlar al alumno, por supuesto, pero permitiendo que rinda tranquilo, como si estuviera en una instancia presencial. Si estás nervioso, que sea por el examen y no por un factor de la plataforma. La otra manera de controlar que tiene esta herramienta es a través del registro absoluto de todo: una cámara frontal, una cámara lateral (que es la cámara del celular) y, además, te graba la pantalla. Te detecta la cantidad de dispositivos conectados a la computadora, si hay dos monitores y todas esas ingenierías que arman a veces. Vos hacés el examen y, suponiendo que alguien lo respondiera todo con ChatGPT, al terminar se manda al profesor. Luego, mediante Inteligencia Artificial, el sistema genera una revisión previa de incidencias. El docente revisa todos los resultados y valida cuáles son incidencias reales y cuáles son falsos positivos. ¿Qué es un falso positivo? Por ahí la herramienta detectó que pasó una persona por atrás, pero cuando ves la grabación es un segundo y no pasa nada. Yo siento que la curaduría humana del resultado tiene que estar siempre, porque si no el control queda supeditado ciegamente a la tecnología y a veces es bastante injusto por los imponderables técnicos. Así funciona el sistema que implementamos. Lo importante y representativo es que tenemos mucho volumen: alrededor de 15,000 chicos cursan materias \textit{online}, pero este sistema lo aplicamos específicamente para 1,200 alumnos, que son los de carreras 100\% a distancia. Con lo cual, cuento con una muestra recontra representativa. Ha salido bien y asegura una calidad académica muy similar, o igual, a la presencial. Es muy raro poder escaparse de este control.
\item[\textbf{Entrevistador:}] Claro. Una consulta: en el modo de expulsión, antes de que el alumno sea expulsado si se detectan anomalías a la hora de hacer el examen, ¿recibe alguna advertencia previa o es directamente expulsado?
\item[\textbf{Entrevistado:}] Sí, en el modo de expulsión, cada vez que detecta algo, el sistema te avisa en la pantalla.
\item[\textbf{Entrevistador:}] Excelente. ¿Y de esta herramienta que comentaste, me dijiste que la adoptaron hace cuánto?
\item[\textbf{Entrevistado:}] La implementamos en marzo, hace poquito, estamos en una etapa reciente.
\item[\textbf{Entrevistador:}] OK, o sea que están probando todavía. ¿Qué cosas dirías que les sirvieron más y cuáles te parece que son un desafío o quedan sin resolver en la parte de evaluaciones y de asegurar la integridad?
\item[\textbf{Entrevistado:}] Mira, yo creo que está resuelto en un 90\%. Lo que queda pendiente ahora es formar al docente para que interprete correctamente los resultados. Nos falta esa curva de aprendizaje; hoy estamos muy metidos nosotros ayudando a los docentes a revisar los informes. Lo que falta es dar el salto para que el docente logre más autonomía, pueda ver los resultados solo y sepa separar los casos graves de los falsos positivos. La implementación de todo sistema es muy compleja. Al principio hay una resistencia por parte de los alumnos, a veces ni siquiera malintencionada. Hay alumnos que no quieren que el sistema tome el control o registre su computadora por un tema de protección de datos personales. Otros tienen notebooks del trabajo y no las pueden usar porque tienen restricciones para instalar software. Eso también lo tendrían que tomar en cuenta para la tesis: el hecho de que deban rendir con otro dispositivo y tengan que conseguirse uno. Nosotros entendemos que si uno elige una carrera a distancia, es básico contar con los dispositivos acordes para poder cursar. Tiene esas limitaciones o complejidades a la hora de la gestión del cambio, tanto con los alumnos como con los profesores para capacitarlos a todos. Por eso a veces es más compleja la gestión del cambio que la implementación tecnológica en sí, porque la herramienta funciona, es estable y es simple.
\item[\textbf{Entrevistador:}] Claro. Mencionaste algo que me llamó bastante la atención del sistema: el seguimiento con una cámara por un lado y otra por el otro. ¿Qué casos de copia te parece que cubrió bien? ¿Hay algún tipo de caso que no cubra bien? Estoy pensando, por ejemplo, en el caso de ChatGPT; ¿se ha enganchado a alguien que se haya copiado de esta forma?, ¿cómo lo detectaron?
\item[\textbf{Entrevistado:}] De todas las formas posibles, porque te graba la pantalla completa de forma continua. El problema que tienen la mayoría de las otras herramientas del mercado es que no cuentan con una cámara lateral y no realizan un paneo del ambiente. Te pueden estar filmando la pantalla y a la persona de frente, pero el alumno puede tener un monitor atrás o un teclado al costado y nadie se da cuenta. El bache de la mayoría de las herramientas es ese: no te dan un paneo del escritorio, de las manos de la persona ni del entorno. Creo que el gran diferencial de esta herramienta es ese y, a priori, me parece que no le falta nada.
\item[\textbf{Entrevistador:}] Súper interesante, buenísimo. No sabía que tenían algo así en UADE, está bueno escucharlo de primera mano.
\item[\textbf{Entrevistado:}] Sí, es nuevito. Incluso ---esto de manera interna para la tesis--- es una muy buena herramienta para implementar en las aulas presenciales y poder evaluar con computadoras. Viste que a veces, por ejemplo si tienen que desarrollar código, los profesores se los tenemos que tomar a papel porque si les dejás usar la computadora te lo resuelven con Claude o con cualquier Inteligencia Artificial. Estas herramientas también podrían servir para evaluar con computadora en el aula y, en algún punto, aggiornarse más tecnológicamente a la evaluación.
\item[\textbf{Entrevistador:}] Bien, sí, está bueno. Y pensando en ampliar la evaluación con computadora a personas que están presenciales, ¿cuál sería el beneficio principal? ¿Poder tomar otro tipo de evaluaciones prácticas?
\item[\textbf{Entrevistado:}] Claro, poder tomar evaluaciones más reales. En informática, que puedas codear en una librería específica sin copiarte de una IA, o que puedas usar la IA si el profesor lo autoriza, pero permitiendo que el docente vea exactamente en qué la usás y en qué no. Por ejemplo, yendo a un área temática que domino más como las ciencias económicas (contabilidad o finanzas), se usa mucho Excel y está buenísimo que se pueda usar en la evaluación. A veces no terminamos usándolo en los exámenes presenciales porque si abren Excel el alumno puede tener todo el examen precargado ahí. A eso voy, te permitiría flexibilizar y enriquecer más las evaluaciones.
\item[\textbf{Entrevistador:}] Súper claro. La verdad es que aceleraste un montón las preguntas que teníamos porque respondiste con mucha información cada una de ellas, fue súper amplio. Te contamos un poco más de nuestro lado: nosotros estamos trabajando en un sistema que comparte principios con el \textit{proctoring}, como se llama esta plataforma que utilizan ustedes. Una tecnología que nos llamó mucho la atención es la biometría conductual, que sirve para diferenciar los patrones intrínsecos de las personas a través de la conducta. En relación con lo que comentabas de la resistencia del alumno o este estrés agregado por el exceso de control, estamos trabajando en asegurar la integridad académica y detectar la copia a través de herramientas no invasivas. Por ejemplo, detectar los patrones de uso del teclado, la repetición de las teclas, el movimiento del mouse o el cambio de pestañas; combinando estas señales. Es decir, si hubo un cambio de pestaña, cuánto tiempo duró, si se copió y pegó texto, qué cantidad, si se modificó ese texto o si se dejó tal cual. En vez de contar con una cámara, ir hacia un análisis que sea más anónimo y que resguarde más la privacidad de los datos. Aunque entendemos perfectamente tu punto: para colocar un cero o invalidar un examen, las pruebas tienen que ser contundentes.
\item[\textbf{Entrevistado:}] Sí, ahí el problema ---y te lo digo desde la experiencia--- es qué pasa cuando detectás una anomalía. Supongamos que el sistema te dice: ``Che, Martina, vos cambiaste de pestaña cinco veces''. Si no hay una evidencia visual o real, es muy complejo. En UADE, toda evidencia comprobada de copia implica un cero y recursar la materia (tanto en la modalidad \textit{online} como en general). ¿Cómo le decís a un alumno que tiene un cero y recursa la materia si no tenés la evidencia real donde se muestre su pantalla abriendo ChatGPT, un resumen oculto, que rindió con un compañero o que rindió otra persona por él? Está bueno el enfoque de ustedes, pero no dejen de tener en cuenta qué pasa con la evidencia que genera el sistema a la hora de tener el ida y vuelta con el alumno, porque lo primero que hace el estudiante es apelar y preguntar qué pruebas tenés para acusarlo. Hoy en día, con los casos que hemos tenido, les muestro la grabación de la pantalla y las cámaras, y ante esa evidencia no hay muchas vueltas que dar. El alumno podrá patalear o arrepentirse, pero la prueba está. Sería muy problemático tomar la decisión de invalidar un examen sin tener una evidencia gráfica irrefutable; les podría generar un problema grande, incluso de índole legal. Les sugiero que lo hablen con su tutor para tenerlo en el radar de la tesis.
\item[\textbf{Entrevistador:}] Claro, es muy interesante para tener en cuenta, súper valioso. En nuestro caso, la idea es que la evidencia se presente a través de un \textit{dashboard} con estadísticas de los patrones recabados para que el docente haga una evaluación desde el lado humano y decida si lo considera evaluable o no. Es un híbrido que busca mitigar el desgaste del estudiante que detectamos en las encuestas, buscando algo eficiente pero con más privacidad. De todos modos, queda clarísimo el punto sobre la contundencia de las pruebas en la vida real.
\item[\textbf{Entrevistado:}] Bueno, excelente. Que tengan mucha suerte con la tesis y cualquier cosa me escriben.
\item[\textbf{Entrevistador:}] Te súper agradecemos el tiempo, fue muy interesante conocer más sobre el funcionamiento de UADE Online. Te mandaremos el trabajo cuando lo presentemos, te va a llegar el agradecimiento correspondiente para que puedas ver la perspectiva final desde tu rol. Muchas gracias, un placer.
\item[\textbf{Entrevistado:}] No, por favor, gracias a ustedes. Que les vaya bien, un beso, saludos.
\end{description}

\section*{ENT-02 --- Joaquín Merlo}
\addcontentsline{toc}{section}{ENT-02 --- Joaquín Merlo}
\begin{ficha}{Ficha de entrevista}
\small
\textbf{Código:} ENT-02\\[2pt]
\textbf{Participante:} Joaquín Merlo\\[2pt]
\textbf{Rol / perfil:} docente universitario e investigador en ciencias biomédicas, con formación en Biotecnología y Bioinformática\\[2pt]
\textbf{Fecha:} 10 de junio de 2026\\[2pt]
\textbf{Modalidad:} virtual\\[2pt]
\end{ficha}
\subsection*{Transcripción completa}
\begin{description}[leftmargin=0cm, labelsep=0.5cm, style=nextline]
\item[\textbf{Entrevistadores:}] Para empezar, si querés contarnos hace cuánto das clase, qué clases das, si tomás algún tipo de evaluaciones presenciales o remotas, en qué proporción, etcétera. Contanos un poco de tu día a día docente.
\item[\textbf{Entrevistado:}] Yo doy clase hace más o menos diez años, casi diez años. Antes de la pandemia no había mucha virtualidad, no sé cómo decirlo, mucha ``remotez''. Se empezó a implementar mucho después. De hecho, creo que se estaba empezando a pensar en 2018 o 2019, cuando yo entré en la carrera de Administración de Empresas online. Recién empezaba eso y todo lo demás era presencial en la facultad. Después, con todo el tema de la pandemia, se cambió mucho a la modalidad remota, no solo durante el aislamiento, sino también mucho después. En general, como yo me especializo en biotecnología y bioinformática, las evaluaciones suelen ser presenciales, porque además hay mucho trabajo práctico. Entonces no hay mucha alternativa a la presencialidad. Pero durante esos años me tocó dar muchas clases y muchas evaluaciones de manera remota. De hecho, todavía seguimos haciendo algunas clases remotas que tienen que ver con los cambios en los horarios de cursada. Todavía seguimos implementando algunas cuestiones remotas. También lo vivo mucho como estudiante, porque hago muchos cursos y, en general, los cursos más internacionales o globales suelen seguir siendo virtuales.
\item[\textbf{Entrevistadores:}] Y seguís dando alguna materia \textit{online}, ¿no? Yo me acuerdo de una materia con vos que era online.
\item[\textbf{Entrevistado:}] Sí, Teoría General de Sistemas. La di una vez virtual. Fue interesante, pero fue en 2020 y pico. La materia era virtual, pero como tenía alumnos de carreras virtuales y presenciales, los parciales y finales los rendíamos presencialmente.
\item[\textbf{Entrevistadores:}] ¿Todo era presencial, los parciales y los finales?
\item[\textbf{Entrevistado:}] Sí, los exámenes fueron presenciales. Lo virtual era la cursada, una modalidad diferente.
\item[\textbf{Entrevistadores:}] Las clases, claro. ¿Pero participaste o evaluaste con evaluaciones \textit{online}, por ejemplo, en pandemia?
\item[\textbf{Entrevistado:}] Sí, sí.
\item[\textbf{Entrevistadores:}] ¿Y qué tal la experiencia? ¿Qué cosas te parece que tienen de ventaja y de desventaja, tanto del lado docente como también del lado estudiante?
\item[\textbf{Entrevistado:}] Ahí por ahí hay que partir un poco la experiencia entre lo que es pre y post inteligencia artificial. Al principio, el desafío de evaluar algo virtual era mucho más sencillo, porque eran tal vez las mismas preguntas, o por ahí meterle un poco más de interpretación o integración a los temas, para que no fuera tan sencillo simplemente buscar la respuesta en un PowerPoint o en un apunte. Hoy en día, con la inteligencia artificial, eso tiene que ir un paso más adelante. Es más complejo, es más difícil, pero se puede. Lo que sí requiere es más tiempo para sentarte, pensar un poco más la evaluación, pensar los contenidos y el enfoque que les das. Eso es algo que tal vez la inteligencia artificial no puede hacer: entender cuál es el enfoque que vos quisiste dar durante la cursada. No es algo que vas a tener guardado en ChatGPT, por ejemplo. Entonces creo que va por ahí: diseñar bien la cursada y evaluar en función de eso, no tanto en función de simplemente los conocimientos que estuviste dando.
\item[\textbf{Entrevistadores:}] Está bueno. Es interesante la postura de diseñar las evaluaciones para esta época, y no tanto seguir evaluando como cuando no existía la inteligencia artificial y después ver cómo adaptarlo para que no se copien.
\item[\textbf{Entrevistado:}] Sí, porque ni hablar de un \textit{multiple choice}. Es muy difícil hacer un \textit{multiple choice} a prueba de inteligencia artificial. Pero incluso hoy en día, las preguntas a desarrollar las inteligencias artificiales las agarran al toque también. De hecho, a veces uso herramientas como [nombre de herramienta no identificado con certeza], les cargo todas mis clases y armo esqueletos de evaluaciones en función de las clases que doy, para después retomarlas y rearmarlas en contexto. Entonces, si yo puedo hacer eso, cualquier alumno puede cargar lo mismo y hacer las mismas preguntas. Pero si uno se centra en lo que fue dando en clase, en cómo interpretar, cómo pensar, cómo encarar un problema, cómo integrar lo que viste en la clase uno o en la clase ocho, eso es algo único y propio de lo que vos transmitís. Ahí es donde hay que hacer foco en la evaluación. Si no, te quedás corto para evaluar.
\item[\textbf{Entrevistadores:}] Bien. En las evaluaciones \textit{online} que tomaste o que conocés, ¿qué métodos se usaban para mitigar o prevenir que se afecte la integridad académica?
\item[\textbf{Entrevistado:}] Por un lado, vi estrategias como poner timers en las preguntas, de manera que el hecho de querer buscar algo, preguntar o mandar un mensaje a alguien ya te empiece a dificultar la respuesta. Dependiendo del tipo de pregunta, eso te puede ayudar o no como docente. Después hay plataformas que detectan si cambiás de pestaña, por ejemplo. Lo que no sé es si son muy buenas detectando que tenés dos monitores. Hace poquito hicimos en la facultad una evaluación de inglés, \textit{online}, pero para nosotros. Durante toda la evaluación tenías que estar mirando a la cámara y tenías que filmarte desde arriba para que se vieran tus manos en el teclado. Donde corrías un poco la vista al costado del monitor, te rechazaba. Tenía un sistema de detección de atención o algo así. Si mirabas a otro lado o se escuchaba un ruido, te detectaba que quizá estabas infringiendo alguna regla.
\item[\textbf{Entrevistadores:}] ¿Y qué sucedía cuando pasaba eso? ¿Te advertía, se terminaba el examen?
\item[\textbf{Entrevistado:}] Sí, te tiraba una advertencia: ``detectamos que está pasando esto, tené cuidado''. Lo que me pasó particularmente con eso es que yo no miré para otro lado; o tal vez sí miré, pero cuando uno piensa, por ahí tiene el gesto de mirar o desenfocar la mirada. Y te complica, porque por ahí no estás haciendo trampa y te advierte. Si te advertía demasiado, te cancelaba la evaluación ahí nomás. Entiendo que todas estas herramientas incluyen al final una instancia de revisión. Imagino que clasifican las evaluaciones entre ``cero riesgo'', ``potencial riesgo'' y ``ya hizo trampa''. Las de potencial riesgo tal vez pasan a algún revisor manual o a alguna persona que se sienta a mirar tu examen. Igual son exámenes de 20 o 30 minutos; no van a mirar un examen de cuatro horas para chequear si corresponde darlo de baja o no.
\item[\textbf{Entrevistadores:}] Y en ese rol de estudiante, ¿qué te pareció? ¿Te parece que es un método efectivo para medir la integridad o para asegurar que no haya copia?
\item[\textbf{Entrevistado:}] Desde las dos posturas, como estudiante y como docente, creo que es una buena herramienta. Pero como docente te exige y te pone en una posición de compromiso, donde tenés que presionar al alumno de cierta manera para evitar estas situaciones. Y como alumno también te mete presión. Por ahí no estás haciendo nada raro e igual te pone un flag de advertencia. Entonces te mete una presión que no solo va sobre el examen y los contenidos, sino también sobre estar sentado derecho y hacer todo de determinada manera para que no haya ningún problema. Es un poco tricky. Creo que a medida que vayan evolucionando las inteligencias artificiales también va a haber más herramientas para permitir una evaluación confiable, pero que pueda transcurrir de una manera menos ortopédica, más natural. Igual, cuando tomé evaluaciones, uno se da cuenta cuando hay un uso o una deshonestidad académica, ya sea copia o uso de inteligencia artificial. Es muy fácil darse cuenta siendo una persona.
\item[\textbf{Entrevistadores:}] ¿Y qué hacés ahí? Por ejemplo, detectás que hay una posibilidad de copia o de uso de inteligencia artificial y no estaba permitido. ¿Cómo accionás?
\item[\textbf{Entrevistado:}] Esa es la cuestión. En general, dependiendo de si lo permitís o no, yo suelo permitirlo, porque también es algo que ya está instaurado. Entonces intento regular el uso más que prohibirlo. Cuando hacemos cosas de programación, si el alumno me pone un código, una estrategia, un programa o un paquete que no vimos en la cursada, o me pone un código de 50 líneas cuando se podía resolver con dos comandos, me doy cuenta de que pasó por inteligencia artificial. Ahí depende de cómo queda la respuesta. Si la respuesta está bien y está bien interpretada, suma, no todo el puntaje completo, pero suma. Y si va para cualquier lado, que muchas veces pasa, queda mal. Si después de la cursada uno tiene que recurrir a la inteligencia artificial para resolver la pregunta, lo más probable es que no entienda la respuesta ni lo que hizo la inteligencia artificial. Al final, la respuesta va a estar igualmente mal. Entonces eso también se autorregula un poco. En preguntas de desarrollo es más difícil, pero uno se da cuenta cuando empiezan a aparecer palabras, conceptos o términos que no usaría nadie de menos de 45 o 50 años. Es muy obvio. También hay que ver el contexto del alumno. Es distinto si es un parcial, un final o un recuperatorio. Todo eso influye para ver cuánto impacta esa deshonestidad, entre comillas, en la nota. En general lo voy manejando caso a caso.
\item[\textbf{Entrevistadores:}] Interesante. Te pregunto indagando en esto que decías de que uno siempre se da cuenta. Entiendo que también tiene que ver con tu experiencia, con casi diez años de docencia. Entre estos recursos que hablamos, como cámara o sistemas de \textit{proctoring} que controlan las evaluaciones, ¿hay alguno que creas que es muy eficiente para detectar copia? ¿O alguno que creas que no se está aplicando y sería útil?
\item[\textbf{Entrevistado:}] De nuevo, todo depende de lo que plantees. Si el planteo es ``vos y tu conocimiento'', estas aplicaciones que te obligan a tener siempre la ventana abierta y no poder cambiar de pestaña para mí son las más eficientes, porque es como estar sentado con tu hoja en el banco. A eso le sumaría la obligación de tener la cámara prendida y el micrófono activo, para evitar ayudas. Pero también depende de si estás en primer año de la facultad, haciendo un curso de posgrado o en quinto año de Ingeniería Informática. Conseguir a una persona que sepa lo mismo o más que vos y que se pueda sentar al lado a soplarte la respuesta es igual de difícil que responderlo vos. No es tan fácil. También depende de la pregunta. Implementar una clase, hacer un comando o un código para hacer tal cosa es una pregunta. Una pregunta de desarrollo o interpretación de un caso es otra cosa. Pienso en la cantidad de opciones que uno puede evaluar o decisiones que puede tomar a la hora de implementar un código. Ahí conseguir ayuda también es muy difícil. Lo que sí tal vez debería haber en ese caso es algún sistema para detectar inteligencia artificial o uso de inteligencia artificial en el desarrollo del código. Pero de nuevo: si estás en Programación I y querés aprender la lógica de programación, está bien prohibirlo. Si estás llegando al final de la carrera, no sé ustedes si están trabajando o no, pero yo el 70\% de las cosas que hago es pedirle a una inteligencia artificial que me implemente algo. No me voy a poner a reinventar la rueda si se lo puedo dar al chat para que lo implemente. Obviamente, después de toda la experiencia, uno tiene conocimiento para controlar eso y ver que funcione correctamente. Pero programarlo de cero no tiene sentido. Entonces depende de la instancia de formación profesional, del enfoque de la evaluación, y ahí uno tiene que ajustar el uso de cada herramienta.
\item[\textbf{Entrevistadores:}] Sí, termina siendo real. Por lo menos en nuestra experiencia, los dos trabajamos y ya programar desde cero es raro. Después, cuando acá tenés que escribir código en papel, es como: ``mirá, hay cosas de la sintaxis que ya ni sirve acordarse''.
\item[\textbf{Entrevistado:}] Sí, pero al principio de la carrera es otra cosa. Ahí sí, por ahí hay que sentarse y aprenderlo bien. Porque si no, el día de mañana una inteligencia artificial tiene un problema o un error, como las alucinaciones clásicas, y si vos no sabés detectarlo a tiempo antes de que eso se convierta en un problema peor, tampoco es tan útil.
\item[\textbf{Entrevistadores:}] Total. Me parece súper interesante que noto en vos muchísimo más empuje desde el lado del diseño de la evaluación que del lado del control. Me parece un punto de vista interesante. Hablamos con varias personas. Una de ellas nos contó de la implementación de un sistema en UADE, que creo que debe ser el mismo que vos tuviste como docente. Parece que también lo usan los alumnos \textit{online}: el sistema en el que tenés cámara adelante, cámara de costado, cambio de pestaña, grabación de pantalla, todo controlado. Me parece interesante esa postura versus la tuya, que es más de aceptar que existe el uso de estas herramientas, pero preguntarse qué podemos hacer para que la evaluación tenga sentido en un momento donde se usan.
\item[\textbf{Entrevistado:}] Sí, pero también cada situación es diferente, no solo por el momento de desarrollo de la carrera, sino también por el contexto institucional. Una cosa es dictar un curso de posgrado, por ejemplo, y otra muy diferente es dictar una carrera en una universidad que tiene que estar acreditada. A nivel ministerio, leyes y toda esa parte más burocrática, uno tiene que asegurar ciertos mecanismos para acreditar la carrera o poder expedir un título. Entonces creo que también hay una mezcla de eso. Imagino que va a ir evolucionando. Pero estos aplicativos, bien aplicados y en cierto contexto, me parece que están buenos y son útiles para controlar ciertas cosas que se te pueden ir de las manos. De nuevo, en Introducción a la Programación, si querés evaluar el desarrollo de un código, tenés que saber hacerlo y yo te miro, y lo tenés que tener en la hoja, porque si no, no puedo evaluar que hayas aprendido los conocimientos. Ahí sí aplica un poco más de control. Pero para otro tipo de preguntas, otro tipo de materias o materias más avanzadas, tal vez eso es un control excesivo que no tiene sentido.
\item[\textbf{Entrevistadores:}] Nosotros estamos trabajando en esto que veníamos hablando, que específicamente se llama \textit{proctoring}: una forma de gestionar o asegurar la integridad académica, o intentar hacerlo. Queremos trabajar con métodos no tan invasivos, como un intermedio entre cuidar al estudiante y generar evidencia para la institución que lo está usando. Puede ser una universidad, un posgrado o certificaciones internacionales, donde esto también se aplica frecuentemente. Lo que queremos hacer es tomar otros métodos que no caigan solo en cámara o micrófono, sino también en el uso de patrones, como el movimiento del mouse o el tipeo en el teclado, para intentar hacer un perfil de la persona e identificar si es la misma que usó la plataforma durante el curso o en el examen final. También están estas herramientas que mencionamos: cambio de pestaña, cuánto tiempo se fue de la pestaña, si hubo copia y pega de texto, etcétera. Algo importante es que la idea no parte de acusar a un estudiante de si copió o no copió, porque eso escapa al alcance de la herramienta. La idea es generar evidencia y trazabilidad del trabajo hecho: que se sepa qué evidencia hay de que la persona fue efectivamente autora de cierta evaluación o de cierto contenido. Para eso se medirían varias señales un poco más privadas y que cuiden más la integridad de la persona. También se podría aportar una capa extra sobre qué preguntas tienen alguna chance de haber sido generadas o traídas de otra fuente, y dar una recomendación al docente: qué pregunta quizás podría hacerle a esa persona para profundizar ese contenido, entender si fue autora o no y tener más poder de decisión como docente. Esa es la idea. Me gustaría escuchar tu opinión.
\item[\textbf{Entrevistado:}] Estoy pensando. Hay un montón de cosas que se me escapan, pero por ejemplo, ¿un Control C se puede detectar? O sea, una época que levante los botones que toca la persona. Detectar Control C, Control V.
\item[\textbf{Entrevistadores:}] Sí, claro. Más que nada Control V, porque se puede ver lo que recibe esa página web. Se puede ver si los caracteres fueron puestos uno a uno, con una diferencia de dos milisegundos entre cada tecla, o si fueron 120 caracteres al mismo tiempo. Si tenés una página web propia, podés medir casi todo. Podés poner un \textit{script} que mida un montón de cosas de la persona. Pero si es una web de otros, o un LMS de learning, las cosas básicas que se pueden medir sin demasiado esfuerzo son la velocidad de tipeo, si se copió, si se pegó, cómo es el movimiento del mouse en la pantalla, si te vas de una pestaña a otra, etcétera. Son cosas relativamente fáciles de trackear y guardan más la privacidad de la persona. Por ejemplo, no podés ver exactamente qué escribió.
\item[\textbf{Entrevistado:}] Sí, ahí el único problema, o no sé si problema, es que hay que avisarlo antes para que esa persona sepa qué datos se le están tomando, entre comillas. Igual, incluso sin hacerlo, creo que funcionaría como elemento disuasivo. Podrían no desarrollar nada y decir: ``mirá que estamos midiendo'', y no medir nada. A mí me parece que está bueno porque resuelve un montón de complejidades. Una persona por ahí se sienta a rendir un examen y tiene una computadora para toda la familia, y vos le estás pidiendo que ponga la computadora, un celular que lo mire desde arriba y otro que lo mire de atrás. Además, que no haya ruido. Me pasó cuando tuve que rendir ese examen: tenía un bebé de dos meses. ¿Qué hago? ¿Lo meto en una caja para que no se escuche ningún ruido de fondo? Es muy difícil, muy complejo.
\item[\textbf{Entrevistadores:}] Es terrible. Yo recuerdo en la pandemia, cuando había evaluaciones 100\% virtuales, de las primeras. Yo estudiaba en [otra universidad] y me acuerdo de una materia donde nos hacían poner la cámara adelante, un celular al costado y el micrófono prendido. Como decís, tenías que asegurar que tampoco hubiera ruido, porque si no afectaba a toda la gente que estaba haciendo la evaluación.
\item[\textbf{Entrevistado:}] Claro, te pasan cuatro o cinco ambulancias y es un quilombo. No tiene ningún sentido. Pero sí, creo que algo que se pueda levantar directamente del dispositivo tal vez es más fácil de trackear, sobre todo para lo que es uso de inteligencia artificial, copia o deshonestidad académica. Después, lo que tiene que ver más con el entorno, ahí hay que ver qué señales se pueden capturar con video, reconocimiento de hacia dónde va la cabeza y eso. Pero es un poco más complejo.
\item[\textbf{Entrevistadores:}] Suponiendo que usás una plataforma de este tipo, que captura señales y te informa que hay una probabilidad de que no tengas evidencia suficiente para decir que una consigna en particular la produjo ese alumno; por ejemplo, hay cierta probabilidad de copia. Y te puede decir las señales: que hubo un evento de pegado de contenido después de que salió de la pestaña. ¿Qué harías? ¿Cómo lo tomarías?
\item[\textbf{Entrevistado:}] Ahí, de nuevo, si tengo 100 alumnos y 80 tienen este flag, hay que ver qué cintura tenés para volver a tomar el examen, por ejemplo, con aquellos que hayan demostrado alguna señal. Si son pocos, para mí sería directamente una instancia oral, en otra fecha o en ese mismo momento cuando termina el examen. Ahí es un ``pasa o no pasa''. Si vos lo sabés y lo podés defender oralmente, pasa. También depende de si es una sola pregunta de diez. En ese caso, tal vez es preguntar sobre esa pregunta y medir ahí qué pasa. Si es en todo el examen, tal vez es un oral completo. De nuevo, es parecido a lo que decía antes: si esa persona tuvo que hacer todo el examen con inteligencia artificial, lo más probable es que cuando lo tenga que defender oralmente no lo pueda hacer. Entonces creo que ahí sí sirve. Cuando los cursos virtuales son muy grandes, esto puede ser una gran herramienta para separar. Sobre ese subset que no lo resolvió en buena ley, tal vez ahí sí hacés algo un poco más personalizado, en una instancia cara a cara, no necesariamente presencial, sino cara a cara en una reunión.
\item[\textbf{Entrevistadores:}] Claro, de uno a uno. Buenísimo.
\item[\textbf{Entrevistado:}] Y si no quiere, desaprobar.
\item[\textbf{Entrevistadores:}] Sí, está bien. Tampoco te podés negar. Si te dicen un oral, no hay que negarse.
\item[\textbf{Entrevistado:}] No hay mucha más alternativa. Ahí también es cuestión de dejar las reglas claras. Yo hace mucho tiempo, en todos los exámenes que tomo, agrego un párrafo con las reglas, sea presencial, virtual o presentación oral. Un parrafito que diga: ``tenés que cumplir con esto, esto y esto; las condiciones para aprobar son estas; y si no, la consecuencia es esta''. Tiene que estar de antemano, porque después aparece el ``pero eso no lo decía en ningún lado''. Entonces te metés en un lío. Las herramientas son tan buenas como la forma en que las implementás. También tenés que tener margen para decir: ``tomo todo el examen de nuevo porque se copiaron todos''. Si desde la universidad o desde otro lugar te dicen que no se toma de nuevo, ahí empezás a quedarte sin herramientas. Entonces hay que ponerlo en contexto.
\item[\textbf{Entrevistadores:}] Genial, buenísimo. Creo que estamos bien. Muchas gracias por tu tiempo.
\item[\textbf{Entrevistado:}] No, por favor. Éxitos, chicos. Como sugerencia, mantengan al tanto a todos los que vayan entrevistando. Cada tanto, si hay algún avance o cuando presenten la herramienta, contacten, así podemos ver en qué terminó todo esto. Es interesante.
\item[\textbf{Entrevistadores:}] Sí, muchas gracias. La verdad es que los plazos son bastante cortos. Ahora en septiembre u octubre hacemos una entrega que ya es la final para presentarlo en diciembre. Los vamos a ir manteniendo al tanto y te mandamos el trabajo para que veas para qué usamos tu experiencia.
\item[\textbf{Entrevistado:}] ¿Para qué usaron mis datos? Bueno, suerte, chicos. Éxitos. Gracias.
\item[\textbf{Entrevistadores:}] Tus datos, exactamente. Muchas gracias y felicidades por el bebé. Adiós.
\end{description}

\section*{ENT-03 --- Mg. Verónica Fainberg}
\addcontentsline{toc}{section}{ENT-03 --- Mg. Verónica Fainberg}
\begin{ficha}{Ficha de entrevista}
\small
\textbf{Código:} ENT-03\\[2pt]
\textbf{Participante:} Mg. Verónica Fainberg\\[2pt]
\textbf{Rol / perfil:} Departamento de Educación a Distancia FACE-FAIN, UADE Online\\[2pt]
\textbf{Fecha:} 10 de junio de 2026\\[2pt]
\textbf{Modalidad:} virtual\\[2pt]
\end{ficha}
\subsection*{Transcripción completa}
\begin{description}[leftmargin=0cm, labelsep=0.5cm, style=nextline]
\item[\textbf{Entrevistadores:}] Es un sistema informático para lograr más integridad educativa. La idea es poder tener cierta trazabilidad de la producción en exámenes; juntar evidencia de que una persona es quien realizó el trabajo y que es autora de lo que está escrito en un examen. Va para ese lado: es una herramienta informática.
\item[\textbf{Entrevistada:}] Ok, ahí estaba la vinculación. Te hago una pregunta más: ¿cómo definen en el marco de su proyecto la integridad?
\item[\textbf{Entrevistadores:}] Buena pregunta. Nosotros definimos integridad, en este caso, como cuidar que la evaluación esté realizada por un alumno y que cumpla con las pautas que le fueron dadas. Por ejemplo: no utilizar información externa, no consultar con otra persona. Que las pautas para que la evaluación sea tomada en cuenta por una institución sean las que se siguen en la evaluación. Pero estamos súper dispuestos a escuchar tu definición de integridad, quizá sea más correcta.
\item[\textbf{Entrevistada:}] No, no es correcto o incorrecto. Es entenderlo un poco en el marco de la conversación. Aprovecho que esto está siendo grabado para ser también correcta: yo voy a contestar con honestidad, sinceridad y todo lo que pueda ayudarlos. Siempre desde mi persona, desde lo individual. Si bien tengo una trayectoria y un rol dentro de la universidad, no quiero que se asuma que mi opinión o mi postura representa necesariamente la del resto de los directivos o autoridades de UADE.
Entiendo que está muy alineada, porque por algo trabajo acá y me siento identificada con las decisiones que se toman y con las políticas, pero hago ese disclaimer para no poner en boca de otros opiniones que quizá no compartan al 100\%.
\item[\textbf{Entrevistadores:}] Perfecto, súper claro. Si querés hablamos un poco más de contexto. Me gustaría escuchar un poco más sobre vos, sobre el rol que tenés en la facultad y tu experiencia docente.
\item[\textbf{Entrevistada:}] Te cuento un poco de mí. Actualmente me desempeño como directora del Departamento de Educación a Distancia para la Facultad de Ciencias Económicas y la Facultad de Ingeniería. Entiendo que ese es el motivo por el cual te contactaron conmigo.
Previo a eso fui directora del Departamento de Economía y Finanzas durante ocho años. Antes fui directora de carrera de Economía y Finanzas, antes docente funcional y antes docente por hora. Llevo más de veinte años en la institución como docente.
También trabajé en otras instituciones educativas dando clases y en colegios secundarios. Mis primeros pasos en la docencia fueron hace bastante tiempo, pero tengo buenos recuerdos de esa época y de esa modalidad.
Todos los directivos de UADE dictamos clases, y yo sigo dictando tanto en modalidad presencial como online. En los años que tengo de experiencia dicté clases en diferentes modalidades, te diría que prácticamente en todas: [MRI, transcripción incierta], cursos de verano, pandemia, intensivos de invierno. Desde la gestión educativa también trabajé muchos años en un programa de capacitación de [BUA, transcripción incierta].
\item[\textbf{Entrevistadores:}] Buenísimo, súper interesante. Y desde el rol que tenés en educación a distancia, ¿cuáles serían las decisiones que tomás o qué discusiones solés tener en el día a día? ¿De qué va exactamente el rol?
\item[\textbf{Entrevistada:}] Es un rol de gestión educativa. Es hacer que la magia suceda: que las clases se dicten. Tiene que ver con la búsqueda de docentes, la programación y asignación de cursos, el seguimiento de la acción del docente dentro del aula, la identificación de oportunidades de mejora, tanto desde lo pedagógico como desde lo disciplinar.
Hoy, desde \textit{online}, no tenemos quizás tanto el foco en lo disciplinar porque tenemos directores de departamento disciplinar con quienes trabajamos. Cuando era directora del Departamento de Economía y Finanzas tenía un rol más activo en la disciplina. Sin embargo, trabajamos en conjunto con los directores de las distintas facultades para ver cuestiones disciplinares, formas de evaluación y metodologías de enseñanza. Todo lo que hace a que el proceso de enseñanza-aprendizaje llegue a buen puerto.
\item[\textbf{Entrevistadores:}] Ajá, súper interesante. Y cuando mencionás los métodos de evaluación, ¿qué tipo de evaluaciones hay en educación a distancia en UADE?
\item[\textbf{Entrevistada:}] Como criterios o formatos, nosotros tenemos evaluaciones que pueden ser orales, escritas, por evaluación continua o por desarrollo y defensa de un trabajo integrador. Hay materias que se ajustan mejor a cada formato y hay otras limitantes que quizás tienen que ver con el expertise docente o con la dinámica del curso, que hacen que sea más adecuado uno u otro.
Una materia práctica se puede tomar de forma oral, pero lleva mucho más tiempo que una materia teórica, porque ese desarrollo práctico se tiene que dar en el momento de la evaluación. No quiere decir que no se pueda hacer. De hecho, todas las materias se pueden evaluar con distintas modalidades.
Hay que tener en cuenta que los objetivos de aprendizaje de una materia, dentro del plan de estudios y estipulados en el programa analítico, estén alineados con las estrategias de enseñanza que se utilizan en el aula y luego con la metodología de evaluación.
Si tenemos una materia donde se trabaja toda la parte práctica, en el examen no debería evaluarse solo teóricamente cuáles son los principios. Poder, se puede, pero no estaría alineado con el objetivo. Y al revés: no podemos trabajar desde la teoría y clasificaciones durante todas las clases y en el examen decir ``aplicalo a un caso práctico''.
Dentro del rol docente también hay que transmitir conocimiento y experiencia que permitan vincular esos conocimientos con la aplicabilidad que van a tener en la carrera profesional. Pero esto es una carrera universitaria, no una tecnicatura ni un curso corto. También se desarrollan habilidades a lo largo de la carrera: expresión oral, expresión escrita, razonamiento lógico, trabajo en equipo, con todas las dificultades que eso conlleva, y también integridad y honestidad académica. Eso es un valor importante en la educación. Va más allá del contenido específico: aprendí a programar o no aprendí a programar.
\item[\textbf{Entrevistadores:}] Perfecto. Mencionaste esta área de foco: asegurar que las evaluaciones tengan sentido con el contenido de la materia, con los recursos docentes, etcétera. Mi pregunta es: ¿qué otros desafíos grandes enfrenta la educación a distancia, por ejemplo comparándola con la modalidad presencial?
\item[\textbf{Entrevistada:}] Uno de los principales desafíos que tiene la educación a distancia es el \textit{engagement}, cómo hacer para que el alumno esté enganchado y comprometido durante un proceso de enseñanza atravesado por la tecnología.
En la presencialidad uno tiene mucha información y comunicación no verbal que sucede en el aula. Uno entra a un aula y dice: ``hoy están cansados'', o ``les costó llegar''. Puede leer lo que está sucediendo dentro del grupo. Eso, en un encuentro sincrónico o en un intercambio en un foro, es mucho más complejo.
Entonces está el desafío adicional de cómo obtener esa información y cómo transmitir empatía y acompañamiento al alumno cuando no está el cara a cara, la conversación informal mientras empieza la clase, el acercamiento al final de la clase o en un recreo, cuando un alumno te dice: ``profe, ¿podrías volver a explicar esto?''.
En la virtualidad los canales están, no quiere decir que no estén, pero es más difícil que se dé ese acercamiento que en la presencialidad sería casi natural. En presencial, si explico algo y alguien pone cara rara, puedo decir: ``¿no lo entendiste?, ¿querés que lo repita?''. En la virtualidad, aunque hoy con las cámaras prendidas uno tiene más información que antes, sigue siendo complejo. Es un desafío más que una problemática, porque se trabaja y hay soluciones para reducirlo.
Otro gran desafío de la educación a distancia tiene que ver con la responsabilidad y el compromiso autónomo del alumno. La cantidad de contenidos y horas de estudio en presencial y \textit{online} son las mismas. El trabajo del docente, en volumen de trabajo, material y demás, también ofrece al alumno lo mismo. Sin embargo, los canales son distintos.
En presencialidad, yo llego, me siento, estoy, y por estar físicamente en el lugar recibo y hago un montón de cosas. Después tengo responsabilidades por fuera, como estudiar. En una materia \textit{online} tengo que entrar, conectarme, estar en una clase y prestar atención. Si me disperso, no hay alguien que necesariamente me llame la atención. Si quiero mirar otra cosa de fondo y pierdo la conexión, es responsabilidad mía volver a poner el foco.
También tengo que entrar semana a semana a ver videos, leer material, realizar actividades, ver devoluciones y preguntar. No es lo mismo cuando el docente en el aula dice ``¿hay alguna duda?'' que cuando eso está en un foro o en un servicio de mensajería.
Entonces se necesita un rol más activo del estudiante: hacerse cargo de la propia formación y no quedarse únicamente con lo que tiene fácil al alcance de la mano. Si no clickeo y entro a cada material, no lo veo. En cambio, si estoy sentada en un aula, tendría que cerrar los ojos y taparme los oídos para no recibirlo.
Eso implica un compromiso del alumno, cuando sabemos que las vidas personales, profesionales y académicas no son lineales. Hay que generar una rutina y un espacio que no necesariamente se tiene. Si cursás una materia presencial todos los lunes a la noche, sabés que a las 15:30 vas a estar ahí. Si cursás \textit{online}, tenés que ponerte la alarma para conectarte, entrar a la plataforma, hacer una entrega, dejar de estudiar para otra materia o dejar de hacer el trabajo que pidió tu jefe, o decirle que no a una salida con una pareja, un amigo o un familiar porque tenés que hacer un trabajo práctico.
Requiere orden y compromiso. No digo que al alumno no le interese, sino que necesita un ordenamiento que no todos tienen y que no siempre es fácil.
\item[\textbf{Entrevistadores:}] Muy interesante. Me siento representada. Este cuatrimestre estoy haciendo una materia virtual y la dedicación que le puse versus las presenciales tiene mucha diferencia. En una encuesta me llamó la atención que preguntaran cuál fue el grado de dedicación. Me hizo reflexionar que hay una cuota más de esfuerzo que uno tiene que poner para consultar o estar al día, algo que en presencial se da más solo.
\item[\textbf{Entrevistada:}] Sí. Hay muchas cosas que uno como docente menciona en el aula y algunos alumnos las toman y otros no. Al final uno dice: ``tenía razón el profesor''.
Yo les digo a los alumnos que imaginen que están en un aula: que agarren cuaderno, lapicera y tomen nota, porque esa acción de escribir, resolver y preguntar hace que el conocimiento se fije. Si uno simplemente ve la exposición del docente, no digo que sea como nada, pero es mucho menos que en el aula. En el aula no está solo el profesor: están todos tus compañeros anotando, entonces te da cosa y anotás, porque el profesor te está mirando.
En la virtualidad podés estar atento, pero de forma pasiva, y de esa manera el aprendizaje no se va a generar.
\item[\textbf{Entrevistadores:}] Y en particular con las evaluaciones, algo que me llama la atención de UADE es que las carreras \textit{online} tengan exámenes presenciales. Sé que para un grupo de personas que hacen carreras \textit{online} esto sí es 100\% virtual, pero me da intriga: ¿por qué todavía deciden que las evaluaciones sigan siendo presenciales? ¿Qué impide pasar al formato 100\% virtual, si es que eso es un objetivo?
\item[\textbf{Entrevistada:}] No sé si es un objetivo o no. Hoy no es un planteo y no te lo podría decir de forma fehaciente.
Sí te digo que es una decisión que se tomó hace tiempo, teniendo en cuenta que el alumno que viene a rendir examen es un alumno de carrera presencial. Es un alumno que asiste para otras materias de forma asidua a UADE. También es darle al examen un marco diferente al resto de las clases y encuentros.
Tiene que ver, por una parte, con las metodologías que se pueden poner en juego en una evaluación presencial, distintas a las online. Pero también tiene que ver con decir: ``este es el momento del examen, nos sentamos en un aula, hay un ordenamiento, está el docente''. También es un espacio donde el docente y el alumno se ven la cara. Eso genera un vínculo distinto.
No digo que uno sea mejor que otro. Nosotros también tomamos exámenes a alumnos de carreras \textit{full online} a distancia y eso tiene sus pros y contras. Pero el alumno de carrera presencial está acostumbrado a tener ese vínculo con el docente y a verlo. Entonces nos parece valioso que en una instancia de evaluación también pueda estar con el docente que lo acompaña durante la cursada. Es una metodología que funciona.
\item[\textbf{Entrevistadores:}] ¿Y cuáles te parece que son esos pros y contras de la evaluación para quienes hacen carrera a distancia y tienen todas las evaluaciones remotas?
\item[\textbf{Entrevistada:}] Para mí, un punto diferencial interesante es que el alumno tenga la posibilidad de conocer personalmente a su docente. No lo pondría como pro o contra del alumno \textit{online}, porque me parece que son perfiles diferentes.
El alumno que decide hacer una carrera a distancia elige una carrera con ventajas de flexibilidad y ubicación, pero sabe que en lo vincular tiene ciertas restricciones. El alumno que decide hacer una carrera presencial está valorando no solo el conocimiento, sino también las dinámicas sociales que se dan en el aula, el networking, el vínculo cara a cara con el docente. Entonces, en ese marco, esa instancia de evaluación también está presente.
En algún momento viene, ve a sus compañeros y ve a sus docentes. En ese sentido me parece una ventaja.
Hoy, dentro de las carreras \textit{full online}, utilizamos una herramienta de monitoreo en el momento de los exámenes. Entonces, el problema que planteaste al principio sobre integridad u honestidad no sería necesariamente un diferencial entre una modalidad y otra. No sería decir: ``presencialmente, como están rindiendo al lado, no se copian; y desde sus casas sí''. Existen otras herramientas que nos permiten ver qué está haciendo el alumno durante el examen.
El diferencial, para mí, va más por el espacio, el ordenamiento y el marco que te da rendir dentro del aula universitaria. Si rendís en tu casa, hay alumnos con realidades diferentes: a veces está la familia, la pareja, los padres, el hermano, entra alguien, sale alguien, te desconcentrás, escuchás una televisión. Eso sucede porque estás en tu casa. Entonces tenés que poner cierto marco y decir: ``estoy rindiendo un examen, por favor no me interrumpan, cierro la puerta, estoy en un lugar tranquilo, tengo lo que necesito para rendir''. El aula ya te da ese marco.
\item[\textbf{Entrevistadores:}] Nos comentaron que están usando una herramienta para detectar algunos eventos de copia. No que directamente acuse, pero sí que le dé al docente más herramientas. Me dijeron que tiene cámara. ¿Tienen pruebas suficientes para decir si una persona se copió o no? ¿Cuál es tu perspectiva sobre esta herramienta? ¿Qué pros y contras tiene? ¿Qué valorás de la herramienta y qué cosas permanecen no resueltas respecto de la trazabilidad de quién realiza la evaluación?
\item[\textbf{Entrevistada:}] Me voy a parar desde otro lugar para contestarte, por mi sesgo profesional de economista. Los agentes tomamos decisiones maximizadoras en base a incentivos. Hay principios que nos ordenan y nos hacen tomar decisiones, pero también hay contextos de incentivos.
Si estoy rindiendo un examen de una materia teórica o numérica y estoy sola, quizás tengo más incentivos para abrir un libro que si estoy en un aula y hay un docente mirándome. Eso no me hace ni mejor ni peor persona: es un contexto que tiene que ver con las decisiones.
También hay grises. Si estoy rindiendo un examen con una computadora, ¿puedo tener un archivo abierto o no? Cuando un docente diseña un instrumento de evaluación, ese instrumento tiene un objetivo alineado con el aprendizaje del alumno y con su formación. Que utilice o no una herramienta complementaria cambia el objetivo del examen.
Puedo diseñar un examen donde se utilice calculadora y otro donde no. Un examen donde se utilice inteligencia artificial y otro donde no. Ahora, si el examen está pensado para que resuelvas a mano sin calculadora, o para que lo hagas con el conocimiento adquirido sin inteligencia artificial, entonces usar esas herramientas hace que el instrumento ya no sirva para ver si efectivamente se generó el aprendizaje buscado.
Tener una herramienta de monitoreo nos permite darle un marco certero y eliminar ambigüedad respecto del examen: esto se puede hacer y esto no se puede hacer. También permite ser claros con el alumno: ``si hacés esto, no estás cumpliendo con el código de conducta de la universidad; es una acción enmarcada dentro de la deshonestidad académica. Si hacés esto otro, está bien porque es lo que se busca con este examen''.
Por eso no me pararía en la postura de evitar que hagan trampa porque está mal o porque no deberían usar esas herramientas. Lo pensaría más desde una formación integral: qué busco lograr, cómo lo voy a lograr, y cómo veo que el alumno efectivamente esté haciendo lo que se preparó para hacer.
No estudiás igual para un examen si sabés que vas a tener una hoja de fórmulas, si vas a tener apuntes al lado o si vas a tener inteligencia artificial. Y no todos los alumnos piensan igual ni responden a los incentivos de la misma forma.
Entonces, tener una herramienta de monitoreo nos permite enmarcar el examen y asegurarnos de que el alumno está haciendo lo que tiene que hacer. También permite equidad horizontal: que todos los alumnos rindan con el mismo material, en las mismas condiciones, y obtengan la nota en función de esa preparación.
\item[\textbf{Entrevistadores:}] Buenísimo. ¿Cómo es el output que les da esta herramienta? ¿Les sirve para tomar decisiones? Entiendo el objetivo que comentabas: asegurar que el alumno esté siendo evaluado de la forma en que fue pensado el examen. ¿Qué recursos les proporciona a los docentes? ¿Y cuál fue la respuesta de los docentes ante esta herramienta? ¿Les es útil?
\item[\textbf{Entrevistada:}] La verdad que sí, la devolución de los docentes fue muy buena. Les gustó esta incorporación. Hay un tema ético que está presente. Un docente hace mucho esfuerzo armando una materia, enseñando y evaluando. Cuando nota que un alumno fue deshonesto de alguna manera, siente una decepción a nivel profesional.
Uno está formando un colega que va a tener un título en la universidad en la que uno trabaja. Si realiza un acto de deshonestidad académica, está traicionando principios que nos acompañan. Entonces, cuando al docente se le dice ``te estoy dando una herramienta que te ayuda a poner un marco a una evaluación'', la incorporación y aceptación fue muy buena.
Están contentos con tener esta alternativa. También les permite tomar otro tipo de evaluaciones. Si el alumno tiene acceso a todos los recursos en la computadora, yo tengo que suponer que los va a usar. Entonces tengo que pensar una evaluación donde los utilice.
Si quiero que desarrolle otro tipo de habilidad o conocimiento, o chequear que incorporó una técnica sin apoyo adicional, necesito asegurarme de que no la utilice. Y ahora puedo hacerlo.
\item[\textbf{Entrevistadores:}] Muy interesante, muy completas tus respuestas. De mi lado ya no tengo más preguntas a nivel entrevista. Si tenés cinco minutos más, te cuento un poco sobre nuestra solución y me das tu opinión. Si no, lo dejamos acá.
\item[\textbf{Entrevistada:}] Dame un segundo, me parece que necesitan la sala. Si no, me corro a otro lugar.
\item[\textbf{Entrevistadores:}] Dale, perfecto.
\item[\textbf{Entrevistada:}] Dame un segundo que me muevo de sala y ahora retomamos.
\item[\textbf{Entrevistadores:}] Perfecto, gracias.
\item[\textbf{Entrevistada:}] Ya estoy de vuelta. Disculpá la demora. Tenía reservada una sala de reuniones y volví a la oficina.
\item[\textbf{Entrevistadores:}] No te preocupes. Te cuento nuestra solución y vos me darás tu opinión. Todavía no tenemos todo 100\% definido. Estamos en una etapa temprana, definiendo lineamientos que van a orientar la parte técnica posterior. Por ejemplo, este sábado tenemos una entrega donde estamos definiendo un marco teórico y todavía no vamos a exponer los resultados del user research.
La idea general es detectar autoría en evaluaciones \textit{online}, no de la misma forma que el sistema implementado en UADE, con cámara o bloqueo del navegador, sino de una forma más pasiva o menos invasiva.
Tiene que ver con biometría conductual. Por un lado, generar un perfil del alumno según cómo tipea, los movimientos del teclado, cómo mueve el mouse, si sale o entra de una pestaña, si copia y pega texto, etcétera.
Sería información menos accionable que la que ustedes tienen hoy, porque no hay un video mostrando que alguien abrió ChatGPT. Pero ayudaría al docente a tener más contexto sobre cómo se resolvió cada consigna en un examen.
El público no sería exactamente algo como UADE, porque ya implementaron una solución distinta, que usa cámara y requiere revisión de video por parte de un docente. Eso no todas las instituciones están dispuestas a aplicarlo y también requiere tiempo docente.
Quizás aplicaría a certificaciones \textit{online} o evaluaciones más masivas, donde hay más estudiantes y se necesitan reportes más generales, sin revisar un video por alumno. La idea sería tener \textit{insights} del trayecto para llegar a una respuesta en un examen. A partir de eso, un docente podría elegir complementar con una pregunta extra, hacer una instancia oral, etcétera.
También se podrían obtener estadísticas de todo el curso para tener una noción general de cómo fue el examen: por ejemplo, cuántas personas parecería que usaron otra herramienta, y entender si fue un problema generalizado o algo realizado por una sola persona.
Me gustaría escuchar tu opinión: si te parece lógico, si tiene algún beneficio frente a otras soluciones más completas como la que usan ustedes, o si tenés preguntas.
\item[\textbf{Entrevistada:}] La información siempre es valiosa. Entre no tenerla y tenerla, siempre suma. Pero la pregunta que viene después es: ¿para qué voy a usar esa información y para qué puedo usarla? ¿Para qué voy a gastar recursos en obtener información que después no voy a poder utilizar?
Si me pongo en el lugar de profesora y me decís: ``hay un 30\% de probabilidad de que este alumno haya [fragmento inaudible]'', ¿qué puedo hacer con eso? ¿Puedo anular el examen? ¿Puedo invalidarle el punto? Si no puedo hacer nada, quizás prefiero no saberlo.
\item[\textbf{Entrevistadores:}] Es muy bueno el punto. Ayer lo discutíamos con mi compañero de equipo. Para nosotros el resultado de esto no sería decir ``anulo o no anulo un examen''. Nos parece que no es suficiente para poner un cero, como cuando en un aula ves a una persona sacando un celular y podés anular el examen.
Esto tendría que ver más con una instancia complementaria: recomendar una pregunta extra en el momento, como ``¿por qué elegiste esto y no esto?'', relacionada con el material. O, en otros casos, una instancia oral. Por ejemplo, si en una pregunta aparece la posibilidad de que la persona haya pegado contenido de internet, hacerle preguntas para triangular si conoce o no el contenido y, a partir de eso, tomar una decisión.
Pero quedaría del lado del docente. La herramienta no diría automáticamente ``no hace el examen'' ni tomaría esa decisión.
\item[\textbf{Entrevistada:}] Claro. Ahí también me pongo en rol de directiva. Los docentes pueden tener criterios diferentes. Si no hay una indicación precisa de decir ``si aparece este indicador, tenés que actuar de esta forma'', quizás un profesor dice ``no importa'' y otro dice ``esto es un error y hay que replantearlo''.
Me parece interesante pensarlo más en exámenes adaptativos. Hoy nosotros usamos [plataforma adaptativa de matemática], que es una plataforma de inteligencia artificial para la enseñanza de matemática. Te va poniendo mayor cantidad de preguntas y ejercicios en función de los resultados que obtenés.
En una guía de trabajos prácticos tradicional tenés, por ejemplo, cinco ejercicios del primer tema. Quizás te resultan fáciles y los hacés rápido, o quizás te cuestan y con cinco no es suficiente. Entonces la plataforma adapta la dificultad y la cantidad de práctica en función de tu resultado.
Algo que se me ocurre, escuchándote, es que si hay determinados indicadores, automáticamente la plataforma tome la decisión de decir: ``te sumo esta pregunta oral'' o ``te sumo esta pregunta''. Que no quede solo en un criterio subjetivo de interpretación posterior. Te lo digo como idea mientras te escuchaba, para que no quede en ``tengo esta información, ¿qué hago con ella?''. Porque si no, de alguna manera es imprecisa.
\item[\textbf{Entrevistadores:}] Súper interesante. No sabía de esa plataforma. Me parece muy bueno que vaya sugiriendo automáticamente en vez de tener lo mismo para todos.
\item[\textbf{Entrevistada:}] En general, los exámenes nivelatorios de inglés y otros que existen en distintas plataformas te van poniendo, según cómo respondés, preguntas más fáciles o más difíciles. Uno se pregunta: ``¿cómo puede ser que con diez preguntas sepan si mi nivel es C1, B1 o B2?''. Porque en realidad mis diez preguntas no eran iguales que tus diez preguntas.
La primera pregunta puede tener un nivel similar, pero según cómo voy respondiendo, el sistema decide si sigue por ahí, sube la dificultad o baja. Esa podría ser una alternativa: si voy a hacer una implementación que tiene un costo en dinero y tiempo, que me permita después hacer algo con la información que obtengo.
\item[\textbf{Entrevistadores:}] Lo veo. Tomo esa postura de que la herramienta haga algo, en lugar de simplemente generar un reporte posterior y dejar todo en manos del docente. En estas plataformas, ¿las preguntas las diseña el docente o también la inteligencia artificial?
\item[\textbf{Entrevistada:}] No lo sé. No las diseña el docente usuario. Hay un equipo de profesionales del área que trabajó en eso. Si yo entro a la plataforma, las preguntas ya están; no las tengo que cargar.
Fue realizada por un equipo que contenía matemáticos. Lo que no sé, porque tiene que ver con cómo fue diseñada la plataforma, es si se dieron parámetros o se cargaron bancos de preguntas. Te digo: no es el docente usuario, eso seguro. Sí participaron docentes o profesionales que diseñaron o colaboraron con el diseño de la plataforma, pero no sé qué tan explícitas fueron las preguntas que pusieron.
\item[\textbf{Entrevistadores:}] Tengo varias dudas. Por ejemplo, en un examen de desarrollo, hay una consigna donde se detecta que el alumno salió del navegador, volvió, pegó un texto, pasó dos segundos en la pantalla y después pasó a la siguiente pregunta. Es un evento exagerado, pero mostraría alta probabilidad de \textit{copy-paste}.
¿Qué tan cómoda se sentiría la universidad dejando que la inteligencia artificial le repregunte al alumno alguna razón por la cual eligió esa respuesta, o algo que tenga que ver con la producción que hizo en esa pregunta?
\item[\textbf{Entrevistada:}] Como docente universitaria, si vos me decís ``dejo que la plataforma decida'', te digo que no. Que la inteligencia artificial no decida nada, por favor. Ahora, si yo lo puedo parametrizar, sí.
Por ejemplo, decir: ``esta es la pregunta corta; de esta pregunta podés preguntarle cuál fue el criterio, qué pasaría si tal cosa, qué harías con esto, por qué no diste esto''. Ahí sí.
A veces también puede ser un cambio en el formato. Dentro de la plataforma podría pedirse una evaluación oral donde el alumno tenga que grabar algo. Ahí corrés el riesgo de que lea, pero si estás tomando datos biométricos, quizás podés darte cuenta si está leyendo o no. Aunque ahí entrás en un círculo.
Como docente, uno podría darle alternativas a la herramienta: si veo que tenés determinado riesgo en cuanto a las acciones que realizaste, te puedo complementar con alguna de estas preguntas de análisis más complejas.
Otra cosa que se puede hacer, y que en diseño de instrumentos de evaluación se hace, es agregar preguntas de validación. A veces en certificaciones, más que en exámenes universitarios, uno completa muchas preguntas de \textit{multiple choice} y aparece una pregunta que en el fondo ya respondió, pero está formulada de otra manera. No es exactamente igual: antes lo preguntó en positivo y ahora en negativo.
Ese tipo de preguntas se incorporan en un instrumento de evaluación y se corrigen en conjunto. Deberías haber respondido las dos correctamente; si no, se invalida, porque si en una ponés que sí y en otra ponés que no, puede ser que haya sido el azar el que contestó.
También podría aparecer más adelante otra pregunta que retome esos conceptos de otra forma. Si las respuestas no son consistentes entre sí, la pregunta no se toma como válida. A veces el alumno puede desarrollar un ejercicio o problema correctamente y después la pregunta es: ``¿qué decisión tomás?, ¿aceptás o no aceptás?''. Contestar mal eso te invalida todo lo demás.
\item[\textbf{Entrevistadores:}] Súper interesante. Te agradezco mucho tu tiempo. Seguramente usemos bastante del contenido que aprendimos de vos y de tu rol para definir cuál va a ser efectivamente la solución que vamos a desarrollar. Te vamos a compartir algo más robusto para que puedas ver cómo usamos tu experiencia. De mi lado no tengo más preguntas. Si tenés algún comentario extra, buenísimo; si no, muchísimas gracias.
\item[\textbf{Entrevistada:}] No, por nada. A disposición para cualquier duda que tengan. Me gustaría leer los avances y después el trabajo terminado, para ver cómo resultaron todas estas indagaciones.
\item[\textbf{Entrevistadores:}] Súper, buenísimo. Muchísimas gracias, que termines bien.
\item[\textbf{Entrevistada:}] De nada, igualmente.
\end{description}
